%============================================================
% presentation.tex
% Présentation Beamer : Tri Bitonique & Minage Blockchain en Mercury
% Compiler : pdflatex presentation.tex (2 passes)
%============================================================
\documentclass[aspectratio=169,12pt]{beamer}

% ── Thème ───────────────────────────────────────────────────
\usetheme{Madrid}
\usecolortheme{default}
\usefonttheme{professionalfonts}

% ── Packages ────────────────────────────────────────────────
\usepackage[utf8]{inputenc}
\usepackage[T1]{fontenc}
\usepackage[french]{babel}
\usepackage{amsmath,amssymb}
\usepackage{listings}
\usepackage{xcolor}
\usepackage{tikz}
\usepackage{pgfplots}
\usepackage{algorithm}
\usepackage{algpseudocode}
\usepackage{booktabs}
\usepackage{multicol}
\usepackage{fontawesome5}
\usepackage{tcolorbox}
\usepackage{array}

\pgfplotsset{compat=1.18}
\usetikzlibrary{shapes,arrows,positioning,calc,shadows,
                decorations.pathmorphing,fit,backgrounds}
\tcbuselibrary{skins}

% ── Palette de couleurs ─────────────────────────────────────
\definecolor{MercuryBlue}{RGB}{13,71,161}
\definecolor{MercuryCyan}{RGB}{0,131,143}
\definecolor{BlockGreen}{RGB}{27,94,32}
\definecolor{BlockOrange}{RGB}{230,81,0}
\definecolor{LightBlue}{RGB}{227,242,253}
\definecolor{LightGreen}{RGB}{232,245,233}
\definecolor{LightOrange}{RGB}{255,243,224}
\definecolor{CodeBg}{RGB}{245,245,245}
\definecolor{KeywordColor}{RGB}{13,71,161}
\definecolor{CommentColor}{RGB}{56,142,60}
\definecolor{StringColor}{RGB}{183,28,28}
\definecolor{DarkGray}{RGB}{33,33,33}

% ── Couleurs Beamer personnalisées ───────────────────────────
\setbeamercolor{palette primary}{bg=MercuryBlue, fg=white}
\setbeamercolor{palette secondary}{bg=MercuryCyan, fg=white}
\setbeamercolor{palette tertiary}{bg=MercuryBlue!80, fg=white}
\setbeamercolor{palette quaternary}{bg=MercuryBlue, fg=white}
\setbeamercolor{structure}{fg=MercuryBlue}
\setbeamercolor{section in toc}{fg=MercuryBlue}
\setbeamercolor{alerted text}{fg=BlockOrange}
\setbeamercolor{block title}{bg=MercuryBlue, fg=white}
\setbeamercolor{block body}{bg=LightBlue}
\setbeamercolor{block title alerted}{bg=BlockOrange, fg=white}
\setbeamercolor{block body alerted}{bg=LightOrange}
\setbeamercolor{block title example}{bg=BlockGreen, fg=white}
\setbeamercolor{block body example}{bg=LightGreen}
\setbeamercolor{frametitle}{bg=MercuryBlue, fg=white}
\setbeamercolor{title}{fg=white}
\setbeamercolor{subtitle}{fg=MercuryCyan!30}
\setbeamercolor{author}{fg=white}
\setbeamercolor{date}{fg=white!80}
\setbeamercolor{institute}{fg=white!90}
\setbeamercolor{footline}{bg=MercuryBlue!10, fg=MercuryBlue}
\setbeamercolor{headline}{bg=MercuryBlue}

% ── Polices Beamer ──────────────────────────────────────────
\setbeamerfont{title}{size=\Large, series=\bfseries}
\setbeamerfont{subtitle}{size=\normalsize}
\setbeamerfont{frametitle}{size=\large, series=\bfseries}
\setbeamerfont{block title}{size=\normalsize, series=\bfseries}
\setbeamerfont{footline}{size=\tiny}

% ── Pied de page personnalisé ────────────────────────────────
\setbeamertemplate{footline}{
  \leavevmode
  \hbox{
    \begin{beamercolorbox}[wd=0.33\paperwidth,ht=2.5ex,dp=1ex,center]
        {footline}
      \insertshortauthor
    \end{beamercolorbox}
    \begin{beamercolorbox}[wd=0.34\paperwidth,ht=2.5ex,dp=1ex,center]
        {footline}
      \insertshorttitle
    \end{beamercolorbox}
    \begin{beamercolorbox}[wd=0.33\paperwidth,ht=2.5ex,dp=1ex,right]
        {footline}
      \insertframenumber{}/\inserttotalframenumber\hspace*{2ex}
    \end{beamercolorbox}
  }
}
\setbeamertemplate{navigation symbols}{}

% ── Bullets ─────────────────────────────────────────────────
\setbeamertemplate{itemize item}{\color{MercuryBlue}\faAngleRight}
\setbeamertemplate{itemize subitem}{\color{MercuryCyan}\faAngleDoubleRight}
\setbeamertemplate{enumerate item}{\color{MercuryBlue}\textbf{\insertenumlabel.}}

% ── Style listings Mercury ───────────────────────────────────
\lstdefinelanguage{Mercury}{
  morekeywords={module,interface,implementation,import_module,
                pred,func,type,inst,mode,det,semidet,nondet,
                multi,is,if,then,else,true,fail,not,
                in,out,di,uo,int,float,char,string,bool,
                list,io,yes,no},
  sensitive=true,
  morecomment=[l]{\%},
  morestring=[b]"
}

\lstset{
  language=Mercury,
  backgroundcolor=\color{CodeBg},
  basicstyle=\ttfamily\scriptsize,
  keywordstyle=\color{KeywordColor}\bfseries,
  commentstyle=\color{CommentColor}\itshape,
  stringstyle=\color{StringColor},
  numberstyle=\tiny\color{gray},
  numbers=left, numbersep=5pt,
  frame=single, framerule=1pt,
  rulecolor=\color{MercuryBlue},
  breaklines=true,
  tabsize=2,
  showstringspaces=false,
  xleftmargin=12pt
}

% ── Commandes utiles ────────────────────────────────────────
\newcommand{\mercury}{\textsc{Mercury}}
\newcommand{\Oh}[1]{\mathcal{O}\!\left(#1\right)}
\newcommand{\code}[1]{\texttt{\color{KeywordColor}#1}}

% ── Méta-données ────────────────────────────────────────────
\title[Tri Bitonique \& Blockchain en Mercury]{%
  \faCode\quad Tri Bitonique \& Minage de Blockchain\\
  \large Implémentation experte en \mercury}
\subtitle{INF322 — Programmation Logique et Fonctionnelle}
\author[INF322]{Cours INF322}
\institute[UYI]{Université de Yaoundé I\\Département d'Informatique}
\date{\today}

%============================================================
\begin{document}
%============================================================

%------------------------------------------------------------
% DIAPOSITIVE TITRE
%------------------------------------------------------------
{
\setbeamertemplate{footline}{}
\setbeamertemplate{headline}{}
\begin{frame}[plain]
\begin{tikzpicture}[remember picture,overlay]
  \fill[MercuryBlue] (current page.north west) rectangle
       ([yshift=-0.55\paperheight]current page.north east);
  \fill[MercuryCyan!20] (current page.south west) rectangle
       ([yshift=0.15\paperheight]current page.south east);
  \fill[white,opacity=0.05]
       ([xshift=0.75\paperwidth,yshift=-0.3\paperheight]current page.north west)
       circle (0.4\paperheight);
  \fill[white,opacity=0.04]
       ([xshift=0.1\paperwidth,yshift=-0.15\paperheight]current page.north west)
       circle (0.2\paperheight);
\end{tikzpicture}

\vspace{0.5cm}
\titlepage
\end{frame}
}

%------------------------------------------------------------
% PLAN
%------------------------------------------------------------
\begin{frame}{Plan de la présentation}
  \tableofcontents[hideallsubsections]
\end{frame}

%============================================================
\section{Introduction à Mercury}
%============================================================

\begin{frame}{\faBook\quad Qu'est-ce que Mercury ?}
  \begin{columns}[T]
    \begin{column}{0.52\textwidth}
      \begin{block}{Définition}
        \mercury{} est un langage \textbf{logique-fonctionnel pur}
        développé à l'Université de Melbourne (1993).\\[0.3em]
        Il combine :
        \begin{itemize}
          \item La logique de \textbf{Prolog}
          \item Le typage de \textbf{Haskell}
          \item Les performances de \textbf{C}
        \end{itemize}
      \end{block}
      \vspace{0.3em}
      \begin{alertblock}{Originalité}
        Chaque prédicat déclare son \textbf{mode} et son
        \textbf{déterminisme} — vérifiés à la \emph{compilation}.
      \end{alertblock}
    \end{column}
    \begin{column}{0.44\textwidth}
      \begin{exampleblock}{Structure de base}
\begin{lstlisting}
:- module hello.
:- interface.
:- import_module io.
:- pred main(
    io::di, io::uo) is det.

:- implementation.
main(!IO) :-
  io.write_string(
    "Bonjour Mercury!\n",
    !IO).
\end{lstlisting}
      \end{exampleblock}
    \end{column}
  \end{columns}
\end{frame}

%------------------------------------------------------------
\begin{frame}{\faCogs\quad Concepts fondamentaux}
  \begin{columns}[T]
    \begin{column}{0.48\textwidth}
      \begin{block}{\faCheckCircle\quad Modes}
        \begin{tabular}{ll}
          \code{in}  & paramètre d'entrée \\
          \code{out} & paramètre de sortie \\
          \code{di}  & entrée destructive (IO) \\
          \code{uo}  & sortie unique (IO) \\
        \end{tabular}
      \end{block}
      \vspace{0.4em}
      \begin{block}{\faListOl\quad Déterminisme}
        \begin{tabular}{ll}
          \code{det}    & toujours 1 solution \\
          \code{semidet}& 0 ou 1 solution \\
          \code{nondet} & 0 à N solutions \\
          \code{multi}  & 1 à N solutions \\
        \end{tabular}
      \end{block}
    \end{column}
    \begin{column}{0.48\textwidth}
      \begin{exampleblock}{\faCode\quad Exemple de déclaration}
\begin{lstlisting}
% Predicat deterministe
:- pred tri(
  list(int)::in,
  list(int)::out) is det.

% Predicat semi-deterministe
:- pred trouve(
  int::in,
  list(int)::in) is semidet.

% Fonction pure
:- func double(int) = int.
double(X) = X * 2.
\end{lstlisting}
      \end{exampleblock}
    \end{column}
  \end{columns}
\end{frame}

%============================================================
\section{Tri Bitonique}
%============================================================

\begin{frame}{\faSortAmountDown\quad Qu'est-ce que le tri bitonique ?}
  \begin{block}{Séquence bitonique}
    Une séquence est \textbf{bitonique} si elle croît puis décroît
    (ou peut s'y ramener par rotation) :
    $$a_1 \leq \cdots \leq a_k \geq a_{k+1} \geq \cdots \geq a_n$$
  \end{block}

  \vspace{0.5em}
  \begin{columns}[T]
    \begin{column}{0.48\textwidth}
      \begin{exampleblock}{\faCheckCircle\quad Avantages}
        \begin{itemize}
          \item \textbf{Parallélisable} : réseau de comparateurs
          \item $\Oh{\log^2 n}$ étapes parallèles
          \item Idéal pour \textbf{GPU / FPGA}
          \item Implémentation régulière et prévisible
        \end{itemize}
      \end{exampleblock}
    \end{column}
    \begin{column}{0.48\textwidth}
      \begin{alertblock}{\faExclamationTriangle\quad Limites}
        \begin{itemize}
          \item $\Oh{n\log^2 n}$ séquentiel
          \item Moins efficace que merge sort ($\Oh{n\log n}$)
          \item Taille $n$ doit être une puissance de 2
          \item Non stable
        \end{itemize}
      \end{alertblock}
    \end{column}
  \end{columns}
\end{frame}

%------------------------------------------------------------
\begin{frame}{\faSortAmountDown\quad Algorithme — Phase 1 : Construction}
  \begin{center}
  \begin{tikzpicture}[scale=0.72]
    % Fils horizontaux
    \foreach \y in {0,1,2,3,4,5,6,7} {
      \draw[gray!40, thick] (-0.3,\y) -- (10.5,\y);
    }
    % Valeurs initiales
    \foreach \y/\v in {0/5,1/1,2/7,3/3,4/8,5/2,6/6,7/4} {
      \node[left,font=\scriptsize] at (-0.3,\y) {$\v$};
    }
    % Étape 1 : paires
    \node[above,font=\tiny,MercuryBlue] at (1.2,7.6) {\textbf{Étape 1}};
    \foreach \ya/\yb in {0/1,2/3,4/5,6/7} {
      \draw[MercuryBlue,very thick] (1.2,\ya)--(1.2,\yb);
      \filldraw[MercuryBlue] (1.2,\ya) circle (3.5pt);
      \filldraw[MercuryBlue] (1.2,\yb) circle (3.5pt);
    }
    % Étape 2
    \node[above,font=\tiny,MercuryCyan] at (3,7.6) {\textbf{Étape 2}};
    \foreach \ya/\yb in {0/3,1/2,4/7,5/6} {
      \draw[MercuryCyan,very thick] (3,\ya)--(3,\yb);
      \filldraw[MercuryCyan] (3,\ya) circle (3.5pt);
      \filldraw[MercuryCyan] (3,\yb) circle (3.5pt);
    }
    \foreach \ya/\yb in {0/2,1/3,4/6,5/7} {
      \draw[MercuryCyan!70,very thick] (4.2,\ya)--(4.2,\yb);
      \filldraw[MercuryCyan!70] (4.2,\ya) circle (3.5pt);
      \filldraw[MercuryCyan!70] (4.2,\yb) circle (3.5pt);
    }
    % Étape 3 : fusion
    \node[above,font=\tiny,BlockGreen] at (7.2,7.6) {\textbf{Fusion bitonique}};
    \foreach \ya/\yb in {0/7,1/6,2/5,3/4} {
      \draw[BlockGreen,very thick] (6,\ya)--(6,\yb);
      \filldraw[BlockGreen] (6,\ya) circle (3.5pt);
      \filldraw[BlockGreen] (6,\yb) circle (3.5pt);
    }
    \foreach \ya/\yb in {0/3,1/2,4/7,5/6} {
      \draw[BlockGreen!80,very thick] (7.5,\ya)--(7.5,\yb);
      \filldraw[BlockGreen!80] (7.5,\ya) circle (3.5pt);
      \filldraw[BlockGreen!80] (7.5,\yb) circle (3.5pt);
    }
    \foreach \ya/\yb in {0/1,2/3,4/5,6/7} {
      \draw[BlockGreen!60,very thick] (9,\ya)--(9,\yb);
      \filldraw[BlockGreen!60] (9,\ya) circle (3.5pt);
      \filldraw[BlockGreen!60] (9,\yb) circle (3.5pt);
    }
    % Valeurs finales
    \foreach \y/\v in {0/1,1/2,2/3,3/4,4/5,5/6,6/7,7/8} {
      \node[right,font=\scriptsize\bfseries,BlockGreen] at (10.5,\y) {$\v$};
    }
  \end{tikzpicture}
  \end{center}
  \vspace{-0.3em}
  \begin{center}
    \small Réseau de comparateurs pour $n=8$
    \quad|\quad $\log_2(8)=3$ phases
    \quad|\quad $\frac{8\times3\times4}{2}=48$ comparaisons
  \end{center}
\end{frame}

%------------------------------------------------------------
\begin{frame}[fragile]{\faCode\quad Implémentation Mercury — Cœur}
  \begin{columns}[T]
    \begin{column}{0.50\textwidth}
      \begin{exampleblock}{Tri bitonique principal}
\begin{lstlisting}
:- pred bitonic_sort(
  direction::in,
  list(int)::in,
  list(int)::out) is det.

bitonic_sort(_, [], []).
bitonic_sort(_, [X], [X]).
bitonic_sort(Dir, List, Sorted) :-
  length(List, N),
  Half = N / 2,
  split(List, Half, L, R),
  % Phase 1 : sequence bitonique
  bitonic_sort(ascending, L, SL),
  bitonic_sort(descending,R, SR),
  append(SL, SR, Bitonic),
  % Phase 2 : fusion
  bitonic_merge(Dir, Bitonic,
                Sorted).
\end{lstlisting}
      \end{exampleblock}
    \end{column}
    \begin{column}{0.46\textwidth}
      \begin{exampleblock}{Fusion bitonique}
\begin{lstlisting}
:- pred bitonic_merge(
  direction::in,
  list(int)::in,
  list(int)::out) is det.

bitonic_merge(_, [], []).
bitonic_merge(_, [X], [X]).
bitonic_merge(Dir, List, Res) :-
  length(List, N),
  Half = N / 2,
  split(List, Half, L, R),
  compare_and_swap(
    Dir, L, R, NL, NR),
  bitonic_merge(Dir, NL, ML),
  bitonic_merge(Dir, NR, MR),
  append(ML, MR, Res).
\end{lstlisting}
      \end{exampleblock}
    \end{column}
  \end{columns}
\end{frame}

%------------------------------------------------------------
\begin{frame}{\faChartBar\quad Complexité du tri bitonique}
  \begin{columns}[T]
    \begin{column}{0.52\textwidth}
      \begin{center}
      \begin{tikzpicture}[scale=0.75]
        \begin{axis}[
          xlabel={Taille $n$},
          ylabel={Opérations},
          xmin=0, xmax=260,
          ymin=0, ymax=50000,
          xtick={64,128,256},
          legend pos=north west,
          legend style={font=\tiny},
          grid=major, grid style={gray!20},
          width=7cm, height=6cm,
          tick label style={font=\tiny},
          label style={font=\small},
          title style={font=\small},
          title={Comparaison des complexités}
        ]
          % O(n log^2 n) - Bitonique
          \addplot[MercuryBlue,very thick,domain=2:256,samples=60]
            {x * (ln(x)/ln(2))^2};
          \addlegendentry{$\Oh{n\log^2 n}$ Bitonique}
          % O(n log n) - Merge sort
          \addplot[BlockGreen,thick,dashed,domain=2:256,samples=60]
            {x * ln(x)/ln(2)};
          \addlegendentry{$\Oh{n\log n}$ Fusion}
          % O(n^2) - Bulles
          \addplot[BlockOrange,thick,dotted,domain=2:220,samples=60]
            {x^2/50};
          \addlegendentry{$\Oh{n^2}$ Bulles}
        \end{axis}
      \end{tikzpicture}
      \end{center}
    \end{column}
    \begin{column}{0.44\textwidth}
      \begin{block}{Tableau comparatif}
        \scriptsize
        \begin{tabular}{lcc}
          \toprule
          \textbf{Algo} & \textbf{Séq.} & \textbf{Par.} \\
          \midrule
          Bitonique  & $\Oh{n\log^2 n}$ & $\Oh{\log^2 n}$ \\
          Fusion     & $\Oh{n\log n}$   & $\Oh{\log n}$   \\
          Rapide     & $\Oh{n\log n}$   & $\Oh{n}$        \\
          Bulles     & $\Oh{n^2}$       & $\Oh{n}$        \\
          \bottomrule
        \end{tabular}
      \end{block}
      \vspace{0.4em}
      \begin{alertblock}{Atout majeur}
        \textbf{Parallélisme maximal} :\\
        $\Oh{\log^2 n}$ étapes\\
        $\Rightarrow$ optimal pour GPU/FPGA
      \end{alertblock}
    \end{column}
  \end{columns}
\end{frame}

%============================================================
\section{Minage de Blockchain}
%============================================================

\begin{frame}{\faLink\quad Architecture d'une Blockchain}
  \begin{center}
  \begin{tikzpicture}[
    bloc/.style={rectangle, rounded corners=4pt,
                 minimum width=2.8cm, minimum height=2.2cm,
                 draw=MercuryBlue, fill=LightBlue,
                 line width=1.2pt, align=center, font=\scriptsize},
    arr/.style={->, thick, MercuryBlue}
  ]
    \node[bloc,fill=LightGreen,draw=BlockGreen] (g) at (0,0) {
      \textbf{Bloc \#0}\\
      \textit{Genesis}\\
      Nonce: 0\\
      Hash: \texttt{00ab1f}\\
      Prev: \texttt{000000}
    };
    \node[bloc] (b1) at (4,0) {
      \textbf{Bloc \#1}\\
      \textit{Alice→Bob}\\
      Nonce: 142\\
      Hash: \texttt{00f3a8}\\
      Prev: \texttt{00ab1f}
    };
    \node[bloc] (b2) at (8,0) {
      \textbf{Bloc \#2}\\
      \textit{Bob→Carol}\\
      Nonce: 387\\
      Hash: \texttt{00d7c2}\\
      Prev: \texttt{00f3a8}
    };
    \node[bloc] (b3) at (12,0) {
      \textbf{Bloc \#3}\\
      \textit{Carol→Dave}\\
      Nonce: 256\\
      Hash: \texttt{00e1b9}\\
      Prev: \texttt{00d7c2}
    };
    \draw[arr] (g.east) -- node[above,font=\tiny,MercuryCyan]
      {prev\_hash} (b1.west);
    \draw[arr] (b1.east) -- node[above,font=\tiny,MercuryCyan]
      {prev\_hash} (b2.west);
    \draw[arr] (b2.east) -- node[above,font=\tiny,MercuryCyan]
      {prev\_hash} (b3.west);

    \node[below=0.3cm of b1, font=\scriptsize,
          color=BlockOrange, align=center]
      {Difficulté : hash\\commence par "00"};
  \end{tikzpicture}
  \end{center}

  \vspace{0.2em}
  \begin{columns}[T]
    \begin{column}{0.48\textwidth}
      \begin{block}{Propriété d'immuabilité}
        Modifier un bloc change son hash
        $\Rightarrow$ invalide \textbf{tous} les blocs suivants
      \end{block}
    \end{column}
    \begin{column}{0.48\textwidth}
      \begin{exampleblock}{Proof-of-Work}
        Trouver $n$ (nonce) tel que :
        $$\text{Hash}(\text{index}\|\text{data}\|n\|\text{prev}) \prec \text{cible}$$
      \end{exampleblock}
    \end{column}
  \end{columns}
\end{frame}

%------------------------------------------------------------
\begin{frame}[fragile]{\faCode\quad Types de données en Mercury}
  \begin{columns}[T]
    \begin{column}{0.50\textwidth}
      \begin{exampleblock}{Définition du type bloc}
\begin{lstlisting}
% Type algebrique pour un bloc
:- type block
    --->    block(
        b_index     :: int,
        b_data      :: string,
        b_nonce     :: int,
        b_hash      :: string,
        b_prev_hash :: string
    ).

% Blockchain = liste de blocs
:- type blockchain == list(block).

% Difficulte = nb de zeros requis
:- type difficulty == int.
\end{lstlisting}
      \end{exampleblock}
    \end{column}
    \begin{column}{0.46\textwidth}
      \begin{exampleblock}{Accès aux champs}
\begin{lstlisting}
% Acces via nom du champ (^)
get_hash(B) = B^b_hash.
get_index(B) = B^b_index.

% Creation d'un bloc
nouveau_bloc(Idx, Data,
             Nonce, Hash,
             Prev) =
  block(Idx, Data,
        Nonce, Hash, Prev).

% Bloc genesis
genesis = block(
  0,
  "Genesis Block",
  0,
  hash_genesis,
  "000000000000").
\end{lstlisting}
      \end{exampleblock}
    \end{column}
  \end{columns}
\end{frame}

%------------------------------------------------------------
\begin{frame}[fragile]{\faHammer\quad Algorithme de minage en Mercury}
  \begin{columns}[T]
    \begin{column}{0.50\textwidth}
      \begin{exampleblock}{Minage (recherche du nonce)}
\begin{lstlisting}
:- pred mine_nonce(
  int::in,       % index
  string::in,    % donnees
  string::in,    % prev_hash
  difficulty::in,
  int::in,       % nonce essai
  int::out,      % nonce trouve
  string::out    % hash valide
) is det.

mine_nonce(Idx,Data,Prev,Diff,
           Try, Nonce, Hash) :-
  H = hash_block(Idx,Data,
                 Try, Prev),
  ( meets_difficulty(H, Diff) ->
    Nonce = Try,
    Hash  = H
  ;
    mine_nonce(Idx, Data, Prev,
      Diff, Try+1, Nonce, Hash)
  ).
\end{lstlisting}
      \end{exampleblock}
    \end{column}
    \begin{column}{0.46\textwidth}
      \begin{exampleblock}{Vérification difficulté}
\begin{lstlisting}
:- pred meets_difficulty(
  string::in,
  difficulty::in) is semidet.

meets_difficulty(Hash, Diff) :-
  Diff > 0,
  string.left(Hash, Diff)
    = Prefix,
  all_zeros(Prefix).

:- pred all_zeros(
  string::in) is semidet.

all_zeros(S) :-
  string.to_char_list(S, Cs),
  all_zero_chars(Cs).

all_zero_chars([]).
all_zero_chars(['0'|Rest]) :-
  all_zero_chars(Rest).
\end{lstlisting}
      \end{exampleblock}
    \end{column}
  \end{columns}
\end{frame}

%------------------------------------------------------------
\begin{frame}[fragile]{\faShieldAlt\quad Validation et sécurité}
  \begin{columns}[T]
    \begin{column}{0.50\textwidth}
      \begin{exampleblock}{Validation de la chaîne}
\begin{lstlisting}
:- pred validate_chain(
  blockchain::in,
  bool::out) is det.

validate_chain([], yes).
validate_chain([_], yes).
validate_chain([B1,B2|R], V) :-
  Expected = hash_block(
    B2^b_index,
    B2^b_data,
    B2^b_nonce,
    B2^b_prev_hash),
  ( Expected = B2^b_hash,
    B2^b_hash = B1^b_prev_hash
    ->
    validate_chain([B2|R], V)
  ;
    V = no
  ).
\end{lstlisting}
      \end{exampleblock}
    \end{column}
    \begin{column}{0.46\textwidth}
      \begin{block}{Détection d'attaque}
        \begin{center}
        \begin{tikzpicture}[scale=0.65]
          \node[draw=BlockGreen,fill=LightGreen,rounded corners,
                minimum width=2cm, align=center, font=\tiny] (b1)
            at (0,2) {Bloc \#1\\Hash: 00f3};
          \node[draw=MercuryBlue,fill=LightBlue,rounded corners,
                minimum width=2cm, align=center, font=\tiny] (b2)
            at (3,2) {Bloc \#2\\Prev: 00f3};
          \draw[->,thick,BlockGreen] (b1)--(b2);

          \node[draw=BlockOrange,fill=LightOrange,rounded corners,
                minimum width=2cm, align=center, font=\tiny] (at)
            at (0,0) {Bloc \#1\\(falsifié)\\Hash: ff99};
          \node[draw=red,fill=red!10,rounded corners,
                minimum width=2cm, align=center, font=\tiny] (b2b)
            at (3,0) {Bloc \#2\\Prev: 00f3\\$\neq$ ff99};
          \draw[->,thick,BlockOrange] (at)--(b2b);
          \node[right,font=\tiny,red] at (4.2,0) {\faTimesCircle};

          \node[above,font=\tiny] at (1.5,2.5) {Chaîne valide};
          \node[below,font=\tiny,red] at (1.5,-0.5) {Attaque détectée};
        \end{tikzpicture}
        \end{center}
      \end{block}
      \begin{alertblock}{Résultat Mercury}
        \code{validate\_chain} $\Rightarrow$ \code{no}\\
        L'attaque est \textbf{automatiquement rejetée}
      \end{alertblock}
    \end{column}
  \end{columns}
\end{frame}

%------------------------------------------------------------
\begin{frame}{\faChartBar\quad Analyse de la difficulté PoW}
  \begin{columns}[T]
    \begin{column}{0.50\textwidth}
      \begin{center}
      \begin{tikzpicture}[scale=0.75]
        \begin{axis}[
          xlabel={Difficulté $d$ (nb de zéros)},
          ylabel={Nonces essayés (moy.)},
          ymode=log, log basis y=10,
          xmin=1, xmax=6,
          xtick={1,2,3,4,5,6},
          grid=major, grid style={gray!20},
          width=7.5cm, height=5.5cm,
          tick label style={font=\tiny},
          label style={font=\small},
          title={Coût du minage vs Difficulté},
          title style={font=\small}
        ]
          \addplot[MercuryBlue,very thick,mark=*,
                   mark options={fill=MercuryBlue}]
            coordinates {(1,16)(2,256)(3,4096)(4,65536)(5,1048576)};
          \addplot[BlockOrange,dashed,domain=1:5,samples=10]
            {16^x};
        \end{axis}
      \end{tikzpicture}
      \end{center}
    \end{column}
    \begin{column}{0.46\textwidth}
      \begin{block}{Propriétés PoW}
        Pour une difficulté $d$ :
        \begin{itemize}
          \item \textbf{Minage} : $\Oh{16^d}$ essais
          \item \textbf{Vérification} : $\Oh{1}$ — immédiate
          \item \textbf{Asymétrie} : coûteux à créer, rapide à vérifier
        \end{itemize}
      \end{block}
      \vspace{0.3em}
      \begin{exampleblock}{Bitcoin (réel)}
        \begin{itemize}
          \item Difficulté $\approx$ 20 zéros hex
          \item $\approx 10^{23}$ hashs/bloc
          \item Ajustée toutes les 2016 blocs
        \end{itemize}
      \end{exampleblock}
    \end{column}
  \end{columns}
\end{frame}

%============================================================
\section{Comparaison et Bilan}
%============================================================

\begin{frame}{\faBalanceScale\quad Comparaison des deux algorithmes}
  \begin{center}
  \begin{tabular}{lll}
    \toprule
    \textbf{Critère} & \textbf{Tri Bitonique} & \textbf{Minage Blockchain} \\
    \midrule
    Paradigme Mercury &
      Récursion pure & Types algébriques + récursion \\
    Déterminisme &
      \code{det} strict & \code{det} (terminaison garantie) \\
    Complexité &
      $\Oh{n\log^2 n}$ & $\Oh{16^d}$ en moyenne \\
    Mémoire &
      $\Oh{n}$ & $\Oh{n}$ blocs \\
    Parallélisable &
      \faCheckCircle\color{BlockGreen}~Oui (natif) &
      \faCheckCircle\color{BlockGreen}~Oui (mining pools) \\
    Pureté Mercury &
      Totalement pure & Pure + IO isolée \\
    Usage réel &
      GPU/FPGA/HPC & Bitcoin, Ethereum \\
    \bottomrule
  \end{tabular}
  \end{center}

  \vspace{0.5em}
  \begin{columns}[T]
    \begin{column}{0.48\textwidth}
      \begin{exampleblock}{\faStar\quad Atout Mercury : Tri bitonique}
        Le système de modes garantit que
        \code{split/4} et \code{compare\_and\_swap/5}
        sont \emph{correctement typés et modés} à la compilation.
      \end{exampleblock}
    \end{column}
    \begin{column}{0.48\textwidth}
      \begin{exampleblock}{\faStar\quad Atout Mercury : Blockchain}
        \code{io::di, io::uo} isole parfaitement les effets
        d'affichage de la logique pure de validation et de minage.
      \end{exampleblock}
    \end{column}
  \end{columns}
\end{frame}

%------------------------------------------------------------
\begin{frame}{\faTerminal\quad Guide de compilation}
  \begin{block}{Prérequis}
    \begin{itemize}
      \item \mercury{} 14.01+ installé (\texttt{mmc} dans le PATH)
      \item Fichiers : \texttt{bitonic\_sort.m} et
            \texttt{blockchain\_mining.m}
    \end{itemize}
  \end{block}

  \vspace{0.4em}
  \begin{exampleblock}{Commandes de compilation et d'exécution}
    \begin{tabular}{ll}
      \toprule
      \textbf{Commande} & \textbf{Description} \\
      \midrule
      \texttt{mmc --make bitonic\_sort} &
        Compiler le tri bitonique \\
      \texttt{./bitonic\_sort} &
        Exécuter le tri bitonique \\
      \texttt{mmc --make blockchain\_mining} &
        Compiler le minage \\
      \texttt{./blockchain\_mining} &
        Exécuter le minage \\
      \texttt{mmc -O2 --make bitonic\_sort} &
        Compilation optimisée \\
      \bottomrule
    \end{tabular}
  \end{exampleblock}

  \vspace{0.3em}
  \begin{alertblock}{Note importante}
    La taille $n$ pour le tri bitonique doit être une
    \textbf{puissance de 2} ($n = 2^k$). Pour d'autres tailles,
    compléter avec des valeurs sentinelles ($+\infty$).
  \end{alertblock}
\end{frame}

%------------------------------------------------------------
\begin{frame}{\faFlagCheckered\quad Conclusion}
  \begin{block}{Ce que nous avons réalisé}
    \begin{itemize}
      \item[\faCheckCircle] \textbf{Tri bitonique} complet en Mercury :
        5 cas de test, vérification automatique, $\Oh{\log^2 n}$ parallèle
      \item[\faCheckCircle] \textbf{Minage Blockchain} avec PoW :
        4 blocs minés, validation d'intégrité, simulation d'attaque détectée
      \item[\faCheckCircle] \textbf{Mercury} : typage fort, modes déclarés,
        pureté garantie, séparation IO/logique
    \end{itemize}
  \end{block}

  \vspace{0.4em}
  \begin{columns}[T]
    \begin{column}{0.48\textwidth}
      \begin{exampleblock}{Points forts de Mercury}
        \begin{itemize}
          \item Erreurs détectées à la \emph{compilation}
          \item Code auto-documenté (modes + déterminisme)
          \item Performance proche du C
          \item Récursion terminale optimisée
        \end{itemize}
      \end{exampleblock}
    \end{column}
    \begin{column}{0.48\textwidth}
      \begin{alertblock}{Extensions possibles}
        \begin{itemize}
          \item SHA-256 réel via FFI C
          \item Tri bitonique sur flottants
          \item Réseau P2P de blockchain
          \item Parallélisme avec \code{par\_conj}
        \end{itemize}
      \end{alertblock}
    \end{column}
  \end{columns}

  \vspace{0.6em}
  \begin{center}
    \large\color{MercuryBlue}
    \textbf{\faCode\quad Merci de votre attention !}\\[0.3em]
    \normalsize\color{gray}
    INF322 — Université de Yaoundé I
    \quad$\cdot$\quad
    \mercury{} 14.01+
  \end{center}
\end{frame}

%------------------------------------------------------------
% DIAPOSITIVES ANNEXES
%------------------------------------------------------------
\appendix

\begin{frame}[fragile]{\faTools\quad Annexe : Fonction de hachage simulée}
  \begin{exampleblock}{Hash polynomial (Mercury — sans SHA-256 natif)}
\begin{lstlisting}
% Hash polynomial deterministe sur une chaine
:- func compute_hash(string) = string.

compute_hash(Input) = HashStr :-
    string.to_char_list(Input, Chars),
    poly_hash(Chars, 0, 0, RawHash),
    AbsHash = abs(RawHash) mod 0xFFFFFFFF,
    HashStr = int_to_hex8(AbsHash).

:- pred poly_hash(list(char)::in, int::in,
                  int::in, int::out) is det.
poly_hash([], _Pos, Acc, Acc).
poly_hash([C|Cs], Pos, Acc, Result) :-
    char.to_int(C, Code),
    NewAcc = (Acc * 1000003
              + Code * int.pow(31, Pos mod 8))
             mod 0x7FFFFFFF,
    poly_hash(Cs, Pos + 1, NewAcc, Result).
\end{lstlisting}
  \end{exampleblock}
  \begin{alertblock}{En production}
    Utiliser \textbf{SHA-256} via une interface C (FFI Mercury) :
    \code{:- pragma foreign\_proc("C", sha256(...))}
  \end{alertblock}
\end{frame}

\end{document}
